\documentclass[12pt,a4paper]{article}

% ===== PAQUETES =====
\usepackage[T1]{fontenc}
\usepackage[utf8]{inputenc}
\usepackage[spanish]{babel}
\usepackage{geometry}
\usepackage{xcolor}
\usepackage{float}
\usepackage{algorithm}
\usepackage{algpseudocode}
\usepackage{amsmath}
\usepackage{amssymb}
\usepackage[most]{tcolorbox}
\usepackage{fontawesome5}
\usepackage{needspace}

% ===== MÁRGENES =====
\geometry{left=2.54cm,right=2.54cm,top=2.54cm,bottom=2.54cm}

\definecolor{azulUNEMI}{RGB}{0,31,63}

% ===== COLORES PSEINT =====
\definecolor{pseKeyword}{RGB}{0,0,153}
\definecolor{pseString}{RGB}{0,102,0}
\definecolor{pseNumber}{RGB}{204,0,0}
\definecolor{pseComment}{RGB}{128,128,128}

% ===== COMANDOS DE COLOR =====
\newcommand{\kw}[1]{\textcolor{pseKeyword}{\textbf{#1}}}
\newcommand{\txt}[1]{\textcolor{pseString}{"#1"}}
\newcommand{\num}[1]{\textcolor{pseNumber}{#1}}
\newcommand{\cmt}[1]{\textcolor{pseComment}{\textit{// #1}}}

% ===== OPERADORES =====
\newcommand{\opAsig}{\textcolor{pseKeyword}{$=$}}
\newcommand{\opMas}{\textcolor{pseKeyword}{$+$}}
\newcommand{\opMenos}{\textcolor{pseKeyword}{$-$}}
\newcommand{\opPor}{\textcolor{pseKeyword}{$*$}}
\newcommand{\opDiv}{\textcolor{pseKeyword}{$/$}}
\newcommand{\opIgual}{\textcolor{pseKeyword}{$=$}}
\newcommand{\opMayor}{\textcolor{pseKeyword}{$>$}}
\newcommand{\opMenor}{\textcolor{pseKeyword}{$<$}}
\newcommand{\opMayorIgual}{\textcolor{pseKeyword}{$\geq$}}
\newcommand{\opMenorIgual}{\textcolor{pseKeyword}{$\leq$}}

% ===== COMANDOS PSEINT =====
\newcommand{\Asignacion}[2]{#1\; \opAsig\; #2}
\newcommand{\Escribir}[1]{\kw{Escribir} #1}
\newcommand{\EscribirSS}[1]{\kw{Escribir} #1 \kw{Sin Saltar}}
\newcommand{\Leer}[1]{\kw{Leer} #1}
\newcommand{\Definir}[2]{\kw{Definir} #1 \kw{Como} \kw{#2}}
\newcommand{\Si}[1]{\kw{Si} #1 \kw{Entonces}}
\newcommand{\Sino}{\kw{Sino}}
\newcommand{\FinSi}{\kw{FinSi}}
\newcommand{\Para}[4]{\kw{Para} #1 \opAsig{} \num{#2} \kw{Hasta} \num{#3} \kw{Con Paso} \num{#4} \kw{Hacer}}
\newcommand{\FinPara}{\kw{FinPara}}
\newcommand{\Repetir}{\kw{Repetir}}
\newcommand{\HastaQue}[1]{\kw{Hasta Que} #1}
\newcommand{\Segun}[1]{\kw{Segun} #1 \kw{Hacer}}
\newcommand{\FinSegun}{\kw{FinSegun}}
\newcommand{\DeOtroModo}{\kw{De Otro Modo:}}
\newcommand{\Subproceso}[1]{\kw{Subproceso} \textbf{#1}}
\newcommand{\FinSubproceso}{\kw{FinSubproceso}}
\newcommand{\Dimension}[1]{\kw{Dimension} #1}

% ===== CONFIGURACIÓN ALGORITHMIC =====
\algrenewcommand\algorithmicindent{1.5em}
\renewcommand{\thealgorithm}{}
\floatname{algorithm}{Algoritmo}

% ===== INDENTACIONES MANUALES =====
\newcommand{\I}{\hspace{1.5em}}
\newcommand{\II}{\hspace{3.0em}}
\newcommand{\III}{\hspace{4.5em}}
\newcommand{\IV}{\hspace{5.5em}}

% ===== ÍNDICE DE FIGURAS APA 7 =====
\newcounter{ejecucion}

\makeatletter
\newcommand{\listadefiguraspseint}{%
	\section*{Índice de figuras}
	\addcontentsline{toc}{section}{Índice de figuras}
	\@starttoc{toe}
}
\makeatother

\newcommand{\InicioEjecucion}[2]{%
	\Needspace{8\baselineskip}
	\refstepcounter{ejecucion}
	\addcontentsline{toe}{section}{Figura \theejecucion. #1}
	\label{#2}
	
	\noindent\textbf{Figura \theejecucion}
	
	\vspace{0.1cm}
	
	\noindent\textit{#1}
	
	\vspace{0.3cm}
}
\newcommand{\consola}[1]{\noindent\textcolor{pseString}{#1}}
\newcommand{\entrada}[1]{\textcolor{black}{#1}}

% ===== MACROS PARA VENTANAS DE EJECUCIÓN PSEINT =====
\newenvironment{PSeIntWindow}{%
	\begin{tcolorbox}[
		breakable,
		enhanced,
		colback=white,
		colframe=black!60,
		boxrule=0.5pt,
		arc=1mm,
		left=0mm,
		right=0mm,
		top=0mm,
		bottom=2mm,
		width=\textwidth
		]
	}{%
	\end{tcolorbox}
}

\newcommand{\BarraPSeInt}{%
	\colorbox{gray!20}{%
		\parbox{\dimexpr\textwidth-2\fboxsep\relax}{%
			\small
			\textcolor{pseString}{\faPlayCircle}
			\hspace{0.3em}
			PSeInt - Ejecutando proceso SISTEMA\_GRUPO3
			\hfill
			\texttt{\_ \hspace{0.5em} \(\square\) \hspace{0.5em} X}
		}
	}
}

\newcommand{\lineaPSeInt}[1]{%
	\begin{flushright}
		\colorbox{gray!20}{\ttfamily\small\textcolor{black}{#1}}
	\end{flushright}
}
%*************************************************

\usepackage{tocloft}

\renewcommand{\cftsecfont}{\color{azulUNEMI}}
\renewcommand{\cftsecpagefont}{\color{azulUNEMI}}

\renewcommand{\cftsubsecfont}{\color{azulUNEMI}}
\renewcommand{\cftsubsecpagefont}{\color{azulUNEMI}}

\usepackage{tocloft}

\renewcommand{\cftsecleader}{\cftdotfill{\cftdotsep}}



%**************************************************
%*************************************************
%***********************************************
\begin{document}
	
	%*******************************************************
	% =====================================================
	% PORTADA EN HOJA BLANCA
	% =====================================================
	
	\begin{titlepage}
		\centering
		
		\vspace*{1cm}
		
		{\Large\bfseries UNIVERSIDAD ESTATAL DE MILAGRO\par}
		\vspace{0.5cm}
		
		{\large\bfseries FACULTAD DE CIENCIAS E INGENIERÍA\par}
		{\large\bfseries CARRERA TECNOLOGÍAS DE LA INFORMACIÓN\par}
		
		\vspace{1.5cm}
		
		{\bfseries TEMA:\par}
		\vspace{0.2cm}
		COMPONENTE PRÁCTICO No. 1
		
		\vspace{1cm}
		
		{\bfseries AUTORES:\par}
		\vspace{0.3cm}
		Albán Coello Diego Andrés\par
		Ayala Pérez Ana Gabriela\par
		Bueno Pincay Solange Marisol\par
		Calero Maldonado Ángel Fernando\par
		Dutan Bonilla Ronald Miguel\par
		González Estrada Ángel Omar (NO TRABAJÓ)\par
		Proaño Alcivar Roberto Celestino\par
		Villega Segura Israel Alfonso (NO TRABAJÓ)\par
		
		\vspace{1cm}
		
		{\bfseries ASIGNATURA:\par}
		Fundamentos de la Programación
		
		\vspace{0.5cm}
		
		{\bfseries DOCENTE:\par}
		Freire Aviles Roger Marcelo
		
		\vspace{0.5cm}
		
		{\bfseries FECHA DE ENTREGA:\par}
		11 de junio 2026
		
		\vspace{0.5cm}
		
		{\bfseries PERÍODO:\par}
		Abril 2026 a Julio 2026
		
		\vfill
		
		{\bfseries MILAGRO -- ECUADOR\par}
		
	\end{titlepage}
	
	
	% =====================================================
	% TABLA DE PARTICIPACIÓN
	% =====================================================
	
	\newpage
	
	\begin{center}
		{\Large\bfseries LINK DE COMPONENTE PRÁCTICO 1}
	\end{center}
	
	\vspace{3cm}
	
	\begin{center}
		{\bfseries Tabla de participación:}
	\end{center}
	
	\vspace{0.5cm}
	
	\begin{center}
		\begin{tabular}{|p{4cm}|p{4cm}|p{3cm}|p{4cm}|}
			\hline
			\textbf{Integrante} & \textbf{Tema} & \textbf{Tiempo} & \textbf{Observación} \\ \hline
			& Introducción & & \\ \hline
			Bueno Solange & & & \\ \hline
			Bueno Solange & & & \\ \hline
			Proaño Roberto & & & \\ \hline
			Ayala Ana & & & \\ \hline
			Albán Diego & & & \\ \hline
			Calero Ángel & & & \\ \hline
			Dután Ronald & & & \\ \hline
			& Conclusión & & \\ \hline
		\end{tabular}
	\end{center}
	
	
	% =====================================================
	% INTRODUCCIÓN Y OBJETIVOS
	% =====================================================
	
	\newpage
	\section{\textcolor{azulUNEMI}{Introducción}}
	
	
	Este proyecto tiene como finalidad integrar los conocimientos adquiridos durante el desarrollo de los Trabajos Prácticos 1 y 2 respectivamente, mediante la construcción de un sistema multitarea en PSeInt que permite ejecutar diferentes ejercicios desde un único menú interactivo.
	
	Para el desarrollo del sistema se aplicaron conceptos fundamentales de programación como el uso de variables, estructuras condicionales, estructuras repetitivas, arreglos unidimensionales y modularización mediante subprocesos. La implementación de estas herramientas permitió organizar el programa de manera eficiente, facilitando la reutilización del código y mejorando su mantenimiento.
	
	El sistema desarrollado incluye cuatro módulos funcionales: conversión de longitudes entre distintas unidades de medida, cálculo de sueldo neto semanal considerando horas extras y descuentos obligatorios, operaciones matemáticas entre vectores y visualización de arreglos en orden inverso. Además, se incorporaron mecanismos de validación de datos para garantizar un funcionamiento adecuado ante posibles errores de ingreso por parte del usuario.
	
	La realización de este proyecto permitió fortalecer las habilidades de análisis, diseño algorítmico y resolución de problemas, competencias esenciales para nuestra formación como futuros profesionales en el área de las Tecnologías de la Información.
	
	\section{\textcolor{azulUNEMI}{Objetivo General}}
	
	
	Desarrollar un sistema multitarea en PSeInt que integre diferentes ejercicios de programación mediante un menú interactivo y subprocesos independientes, aplicando estructuras de control, arreglos y técnicas de modularización para fortalecer los conocimientos fundamentales de programación.
	
		\subsection{\textcolor{azulUNEMI}{Objetivos Específicos}}
	
	
	\begin{itemize}
		\item Diseñar un menú principal que permita acceder a cada una de las funcionalidades del sistema de manera organizada e intuitiva.
		\item Implementar un módulo para la conversión automática de longitudes desde metros hacia pies, pulgadas, yardas y millas.
		\item Desarrollar un módulo de cálculo de nómina semanal considerando horas normales, horas extraordinarias y descuentos obligatorios.
		\item Aplicar el uso de arreglos unidimensionales para realizar operaciones de suma y resta entre vectores.
		\item Utilizar estructuras repetitivas para almacenar, recorrer y mostrar datos en arreglos de tamaño fijo y dinámico.
		\item Emplear subprocesos para modularizar el programa, facilitando la reutilización y el mantenimiento del código.
		\item Incorporar validaciones que permitan controlar errores de entrada y mejorar la confiabilidad del sistema.
	\end{itemize}
	
	
	% =====================================================
	% DESARROLLO DEL SISTEMA
	% =====================================================
	
	\newpage
	
	\section{\textcolor{azulUNEMI}{Desarrollo del sistema}}
	
	\vspace{0.3cm}
	
	% AQUÍ COLOCAS EL CÓDIGO DEL SISTEMA PSEINT FORMATEADO
	
	% AQUÍ VAN TODAS LAS FIGURAS DE EJECUCIÓN
	
	%*********************************************************
	%**********************************************************
	%**********************************************************
	\begin{center}
		{\LARGE\bfseries\textcolor{pseString}{SISTEMA MULTITAREA GRUPO 3}}
	\end{center}
	\vspace{0.5cm}
	\noindent
	\begin{algorithmic}[1]
		
		\State \kw{Algoritmo} Sistema\_Grupo3
		
		\State \I \kw \cmt{Encabezado}
		\State \I \kw \Escribir{\txt{===================================}}
		\State \I \kw \Escribir{\txt{           UNIVERSIDAD ESTATAL DE MILAGRO}}
		\State \I \kw \Escribir{\txt{        FACULTAD DE CIENCIAS DE LA INGENIERIA}}
		\State \I \kw \Escribir{\txt{    INGENIERIA EN TECNOLOGIAS DE LA INFORMACION}}
		\State \I \kw \Escribir{\txt{====================================}}
		\State \I \kw \Escribir{\txt{PROYECTO: COMPONENTE PRACTICO 1}}
		\State \I \kw \Escribir{\txt{GRUPO NRO: 3}}
		\State \I \kw \Escribir{\txt{INTEGRANTES:}}
		\State \I \kw \Escribir{\txt{  DIEGO ANDRES ALBAN COELLO        0959061219}}
		\State \I \kw \Escribir{\txt{  ANA GABRIELA AYALA PEREZ         1719195321}}
		\State \I \kw \Escribir{\txt{  SOLANGE MARISOL BUENO PINCAY     0927255901}}
		\State \I \kw \Escribir{\txt{  ANGEL FERNANDO CALERO MALDONADO  1715227110}}
		\State \I \kw \Escribir{\txt{  RONALD MIGUEL DUTAN BONILLA      0940568066}}
		\State \I \kw \Escribir{\txt{  ROBERTO CELESTINO PROANO ALCIVAR 1206695635}}
		\State \I \kw \Escribir{\txt{  ISRAEL ALFONSO VILLEGA SEGURA    0952182244}}
		\State \I \kw \Escribir{\txt{=======================================}}
		\State \I \kw \Escribir{\txt{SISTEMA INICIALIZADO. Presione una tecla para continuar...}}
		\State \I \kw{Esperar Tecla}
		\State
		
		\State \I \cmt{Declaracion de variables generales}
		\State \I \Definir{opcionSeleccionada}{Entero}
		\State \I \kw \Definir{Salir}{Logico}
		\State \I \Asignacion{Salir}{\kw{Falso}}
		
		
		\State \I \kw \cmt{Ciclo repetitivo hasta que elija Salir}
		\State \I \kw \Repetir
		\State \II \kw{Limpiar Pantalla}
		\State \II \Escribir{\txt{=====================================}}
		\State \II \Escribir{\txt{              MENU PRINCIPAL}}
		\State \II \Escribir{\txt{======================================}}
		\State \II \Escribir{\txt{1. Conversion de Longitudes}}
		\State \II \Escribir{\txt{2. Calcular el sueldo neto semanal de un trabajador}}
		\State \II \Escribir{\txt{3. Suma y resta de dos arreglos unidimensionales}}
		\State \II \Escribir{\txt{4. Inversa de un arreglo de N elementos}}
		\State \II \Escribir{\txt{5. Salir}}
		\State \II \Escribir{\txt{======================================}}
		\State \II \EscribirSS{\txt{Por favor, seleccione una opcion (1-5): }}
		\State \II \Leer{opcionSeleccionada}
		\State
		\State \II \cmt{Llamada a los subprogramas}
		\State \II \Segun{opcionSeleccionada}
		\State \III \num{1}\textbf{:}
		\State \IV SubProgramaConversion
		\State \III \num{2}\textbf{:}
		\State \IV SubProgramaCalculoSueldo
		\State \III \num{3}\textbf{:}
		\State \IV SubProgramaOperacionVectores
		\State \III \num{4}\textbf{:}
		\State \IV SubProgramaVectorInverso
		\State \III \num{5}\textbf{:}
		\State \IV \kw{Limpiar Pantalla}
		\State \IV \Escribir{\txt{Cerrando la sesion con exito de forma segura.}}
		\State \IV \Escribir{\txt{Proyecto desarrollado Por el Grupo 3.}}
		\State \IV \Escribir{\txt{Derechos reservados Grupo 3 febrero-agosto 2026.}}
		\State \IV \Asignacion{Salir}{\kw{Verdadero}}
		\State \III \DeOtroModo
		\State \III \Escribir{\txt{ERROR CRITICO: La opcion [}, opcionSeleccionada, \txt{] no es}}
		\State \III \Escribir{\txt{correcta.}}
		\State \IV \Escribir{\txt{Intente nuevamente introduciendo un numero del 1 a 5.}}
		\State \IV \Escribir{\txt{Presione una tecla para continuar...}}
		\State \IV \kw{Esperar Tecla}
		\State \II \FinSegun
		\State \I \HastaQue{Salir \opIgual{} \kw{Verdadero}}
		\State
		\State \kw{FinAlgoritmo}
		
		
		\State \Subproceso{SubProgramaConversion}
		\State \I \kw{Limpiar Pantalla}
		\State \I \Definir{metrosIngresados}{Real}
		\State \I \Definir{pies, pulgadas, yardas, millas}{Real}
		\State
		\State \I \Escribir{\txt{=====================================}}
		\State \I \Escribir{\txt{                  MODULO 1}}
		\State \I \Escribir{\txt{=====================================}}
		\State \I \EscribirSS{\txt{Ingrese la magnitud en Metros a convertir: }}
		\State \I \Leer{metrosIngresados}
		\State
		\State \I \cmt{Validacion: metros deben ser positivos o 0}
		\State \I \Si{metrosIngresados \opMenor{} \num{0}}
		\State \II \Escribir{\txt{ERROR el numero ingresado debe ser positivo.}}
		\State \I \Sino
		\State \II \cmt{conversion}
		\State \II \Asignacion{pies}{metrosIngresados \opPor{} \num{3.28084}}
		\State \II \Asignacion{pulgadas}{metrosIngresados \opPor{} \num{39.3701}}
		\State \II \Asignacion{yardas}{metrosIngresados \opPor{} \num{1.09361}}
		\State \II \Asignacion{millas}{metrosIngresados \opPor{} \num{0.000621371}}
		\State
		\State \II \Escribir{\txt{-------------------------------------------------}}
		\State \II \Escribir{\txt{RESULTADOS DE LA CONVERSION:}}
		\State \II \Escribir{\txt{  -> Metros origen: }, metrosIngresados, \txt{ METROS}}
		\State \II \Escribir{\txt{  -> Pies         : }, pies, \txt{ PIES}}
		\State \II \Escribir{\txt{  -> Pulgadas     : }, pulgadas, \txt{ PULGADAS}}
		\State \II \Escribir{\txt{  -> Yardas       : }, yardas, \txt{ YARDAS}}
		\State \II \Escribir{\txt{  -> Millas       : }, millas, \txt{ MILLAS}}
		\State \II \Escribir{\txt{======================================}}
		\State \I \FinSi
		\State \I \Escribir{\txt{Presione una tecla para retornar al Menu Principal...}}
		\State \I \kw{Esperar Tecla}
		\State \FinSubproceso
		
		\State \Subproceso{SubProgramaCalculoSueldo}
		\State \I \kw{Limpiar Pantalla}
		\State \I \Definir{totalHorasTrabajadas, costoHoraNormal}{Real}
		\State \I \Definir{sueldoBrutoTotal, deduccionSeguroSocial, sueldoNetoPercibir}{Real}
		\State \I \Definir{horasNormales, horasExtras}{Real}
		\State
		\State \I \Escribir{\txt{==========================================}}
		\State \I \Escribir{\txt{    MODULO 2: LIQUIDACION DE NOMINA SEMANAL}}
		\State \I \Escribir{\txt{===========================================}}
		\State \I \EscribirSS{\txt{Ingrese el total de horas trabajadas en la semana: }}
		\State \I \Leer{totalHorasTrabajadas}
		\State \I \EscribirSS{\txt{Ingrese el costo de la hora laboral estandar (USD): }}
		\State \I \Leer{costoHoraNormal}
		\State
		\State \I \Si{totalHorasTrabajadas \opMenor{} \num{0} \kw{ O } costoHoraNormal \opMenor{} \num{0}}
		\State \II \Escribir{\txt{ERROR: Las horas trabajadas y costos deben ser positivos.}}
		\State \I \Sino
		\State \II \cmt{horas extraordinarias}
		\State \II \Si{totalHorasTrabajadas \opMenorIgual{} \num{40}}
		\State \III \Asignacion{horasNormales}{totalHorasTrabajadas}
		\State \III \Asignacion{horasExtras}{\num{0}}
		\State \III \Asignacion{sueldoBrutoTotal}{horasNormales \opPor{} costoHoraNormal}
		\State \II \Sino
		\State \III \Asignacion{horasNormales}{\num{40}}
		\State \III \Asignacion{horasExtras}{totalHorasTrabajadas \opMenos{} \num{40}}
		\State \III \cmt{recargo horas extras al 150\%}
		\State \III \Asignacion{sueldoBrutoTotal}{(horasNormales \opPor{} costoHoraNormal) \opMas{} (horasExtras \opPor{} costoHoraNormal \opPor{} \num{1.5})}
		\State \II \FinSi
		\State
		\State \II \cmt{Deduccion IESS 9.2\%}
		\State \II \Asignacion{deduccionSeguroSocial}{sueldoBrutoTotal \opPor{} \num{0.092}}
		\State \II \Asignacion{sueldoNetoPercibir}{sueldoBrutoTotal \opMenos{} deduccionSeguroSocial}
		\State
		\State \II \Escribir{\txt{--------------------------------------------------------}}
		\State \II \Escribir{\txt{ROL DE PAGO:}}
		\State \II \Escribir{\txt{  [+] Horas Ordinarias   : }, horasNormales, \txt{ hrs}}
		\State \II \Escribir{\txt{  [+] Horas Extraordinarias: }, horasExtras, \txt{ hrs}}
		\State \II \Escribir{\txt{  [=] Sueldo Bruto        : \$}, sueldoBrutoTotal}
		\State \II \Escribir{\txt{  [-] Descuento IESS 9.2\%: \$}, deduccionSeguroSocial}
		\State \II \Escribir{\txt{   Sueldo Neto Semanal: \$}, sueldoNetoPercibir}
		\State \II \Escribir{\txt{=======================================}}
		\State \I \FinSi
		\State \I \Escribir{\txt{Presione una tecla para retornar al Menu Principal...}}
		\State \I \kw{Esperar Tecla}
		\State \FinSubproceso
		
		\State \Subproceso{SubProgramaOperacionVectores}
		\State \I \kw{Limpiar Pantalla}
		\State \I \cmt{Dimensionamiento de 4 arreglos fijos de 10 celdas}
		\State \I \Dimension{vectorA[\num{10}], vectorB[\num{10}], vectorSuma[\num{10}], vectorResta[\num{10}]}
		\State \I \Definir{vectorA, vectorB, vectorSuma, vectorResta, indiceBucle}{Entero}
		\State
		\State \I \Escribir{\txt{==========================================}}
		\State \I \Escribir{\txt{ MODULO 3: OPERACIONES CON VECTORES (SUMA-RESTA)}}
		\State \I \Escribir{\txt{==========================================}}
		\State
		\State
		\State \I \Escribir{\txt{\textgreater\textgreater\ INGRESAR 10 ELEMENTOS PARA EL VECTOR A:}}
		\State \I \Para{indiceBucle}{1}{10}{1}
		\State \II \EscribirSS{\txt{Ingrese numero VectorA[}, indiceBucle, \txt{]: }}
		\State \II \Leer{vectorA[indiceBucle]}
		\State \I \FinPara
		
		\State
		\State \I \Escribir{\txt{\textgreater\textgreater\ INGRESAR 10 ELEMENTOS PARA EL VECTOR B:}}
		\State \I \Para{indiceBucle}{1}{10}{1}
		\State \II \EscribirSS{\txt{Ingrese numero VectorB[}, indiceBucle, \txt{]: }}
		\State \II \Leer{vectorB[indiceBucle]}
		\State \I \FinPara
		\State \I \cmt{Impresion de resultados}
		\State \I \kw{Limpiar Pantalla}
		\State \I \Escribir{\txt{===========================================}}
		\State \I \Escribir{\txt{        IMPRESION DE VALORES ORDENADOS }}
		\State \I \Escribir{\txt{===========================================}}
		\State \I \Escribir{\txt{Indice | VectorA | VectorB | Suma(A+B) | Resta(A-B)}}
		\State \I \Escribir{\txt{-------------------------------------------------------------------}}
		\State \I \Para{indiceBucle}{1}{10}{1}
		\State \II \Escribir{\txt{[},indiceBucle,\txt{]},vectorA[indiceBucle],\txt{|},vectorB[indiceBucle],\txt{|},
			vectorSuma[indiceBucle],\txt{|},vectorResta[indiceBucle]}
		\State \I \FinPara
		\State \I \Escribir{\txt{===========================================}}
		\State \I \Escribir{\txt{Presione una tecla para retornar al Menu Principal...}}
		\State \I \kw{Esperar Tecla}
		\State \FinSubproceso
		
		\State \Subproceso{SubProgramaVectorInverso}
		\State \I \kw{Limpiar Pantalla}
		\State \I \Definir{dimensionVector}{Entero}
		\State \I \Definir{indice}{Entero}
		\State
		\State \I \Escribir{\txt{==========================================}}
		\State \I \Escribir{\txt{ MODULO4: INVERSO DEL VECTOR INGRESADO}}
		\State \I \Escribir{\txt{==========================================}}
		\State
		\State \I \Repetir
		\State \II \EscribirSS{\txt{Defina el tamano del vector (debe ser positivo): }}
		\State \II \Leer{dimensionVector}
		\State \II \Si{dimensionVector \opMenorIgual{} \num{0}}
		\State \III \Escribir{\txt{ALERTA: dimension N debe ser entero estrictamente positivo.}}
		\State \II \FinSi
		\State \I \HastaQue{dimensionVector \opMayor{} \num{0}}
		\State
		\State \I \Dimension{vector[dimensionVector]}
		\State \I \Definir{vector}{Real}
		\State
		\State \I \Asignacion{indice}{\num{1}}
		\State \I \Repetir
		\State \II \EscribirSS{\txt{Ingrese valor celda dinamica [}, indice, \txt{]: }}
		\State \II \Leer{vector[indice]}
		\State \II \Asignacion{indice}{indice \opMas{} \num{1}}
		\State \I \HastaQue{indice \opMayor{} dimensionVector}
		\State
		\State \I \kw{Limpiar Pantalla}
		\State \I \Escribir{\txt{=========================================}}
		\State \I \Escribir{\txt{VECTOR ORDENADO}}
		\State \I \Escribir{\txt{=========================================}}
		\State \I \Asignacion{indice}{\num{1}}
		\State \I \Repetir
		\State \II \Escribir{\txt{  [*] Elemento Celda [}, indice, \txt{] = }, vector[indice]}
		\State \II \Asignacion{indice}{indice \opMas{} \num{1}}
		\State \I \HastaQue{indice \opMayor{} dimensionVector}
		\State
		\State \I \Escribir{\txt{=========================================}}
		\State \I \Escribir{\txt{ORDEN INVERSO DEL VECTOR}}
		\State \I \Escribir{\txt{=========================================}}
		\State \I \Asignacion{indice}{dimensionVector}
		\State \I \Repetir
		\State \II \Escribir{\txt{  [\#] Indice Invertido [}, indice, \txt{] = }, vector[indice]}
		\State \II \Asignacion{indice}{indice \opMenos{} \num{1}}
		\State \I \HastaQue{indice \opMenor{} \num{1}}
		\State \I \Escribir{\txt{==========================================}}
		\State \I \Escribir{\txt{Presione una tecla para retornar al Menu Principal...}}
		\State \I \kw{Esperar Tecla}
		\State \FinSubproceso
	\end{algorithmic}
	
	
	
	
	%======================================INICIO EJECUCION================		
	
	
	
	
	
	
	\InicioEjecucion
	{EJECUCIÓN INICIAL DEL SISTEMA}
	{ejec:01-ejecucion-inicial-del-sistema}
	
	\begin{PSeIntWindow}
		
		\BarraPSeInt
		
		\vspace{0.3cm}
		
		{\ttfamily
			
			{\color{pseKeyword}*** Ejecución Iniciada. ***}
			
			\vspace{0.3cm}
			
			{\color{pseString}
				=============================================================
				
				\begin{center}
					UNIVERSIDAD ESTATAL DE MILAGRO
					
					FACULTAD DE CIENCIAS DE LA INGENIERÍA
					
					INGENIERÍA EN TECNOLOGÍAS DE LA INFORMACIÓN
				\end{center}
				
				=============================================================
				
				PROYECTO: COMPONENTE PRÁCTICO 1
				
				\hspace*{3cm}GRUPO NRO: 3
				
				INTEGRANTES:
				
				\hspace*{2cm}DIEGO ANDRES ALBAN COELLO \hfill 0959061219
				
				\hspace*{2cm}ANA GABRIELA AYALA PEREZ \hfill 1719195321
				
				\hspace*{2cm}SOLANGE MARISOL BUENO PINCAY \hfill 0927255901
				
				\hspace*{2cm}ÁNGEL FERNANDO CALERO MALDONADO \hfill 1715227110
				
				\hspace*{2cm}RONALD MIGUEL DUTAN BONILLA \hfill 0940568066
				
				\hspace*{2cm}ROBERTO CELESTINO PROAÑO ALCÍVAR \hfill 1206695635
				
				\hspace*{2cm}ISRAEL ALFONSO VILLEGA SEGURA \hfill 0952182244
				
				\vspace{0.3cm}
				
				=============================================================
				
				SISTEMA INICIALIZADO CORRECTAMENTE. Presione una tecla para continuar...
			}
		}
		
	\end{PSeIntWindow}
	%===========================FIN INICIO EJECUCION ==================
	%===========================MENU DE OPCIONES======================
	\vspace{0.2cm}
	\InicioEjecucion
	{MENÚ DE OPCIONES}
	{ejec:02-menu-de-opciones}
	
	\begin{PSeIntWindow}
		
		\BarraPSeInt
		
		\vspace{0.2cm}
		
		\ttfamily
		\small
		
		\consola{========================================================}\par
		
		\begin{center}
			\textcolor{pseString}{MENÚ PRINCIPAL}
		\end{center}
		
		\consola{========================================================}\par
		
		\consola{1. Conversión de Longitudes}\par
		
		\consola{2. Calcular el sueldo neto semanal de un trabajador}\par
		
		\consola{3. Suma y resta de dos arreglos unidimensionales}\par
		
		\consola{4. Inversa de un arreglo de N elementos}\par
		
		\consola{5. Salir}\par
		
		\consola{========================================================}\par
		
		\consola{Por favor, seleccione una opción (1-5): \textgreater}\par
		
		\vspace{2.8cm}
		
		\lineaPSeInt{línea 41 instrucción 1}
		
	\end{PSeIntWindow}	
	%============================MENU DE OPCIONES Opcion 1========================
	\vspace{0.2cm}
	\InicioEjecucion
	{MENÚ DE OPCIONES OPCION 1}
	{ejec:03-menu-de-opciones-opcion-1}
	
	\begin{PSeIntWindow}
		
		\BarraPSeInt
		
		\vspace{0.2cm}
		
		\ttfamily
		\small
		
		\consola{========================================================}\par
		
		\begin{center}
			\textcolor{pseString}{MENÚ PRINCIPAL}
		\end{center}
		
		\consola{========================================================}\par
		
		\vspace{0.2cm}
		
		\consola{1. Conversión de Longitudes}\par
		
		\consola{2. Calcular el sueldo neto semanal de un trabajador}\par
		
		\consola{3. Suma y resta de dos arreglos unidimensionales}\par
		
		\consola{4. Inversa de un arreglo de N elementos}\par
		
		\consola{5. Salir}\par
		
		\vspace{0.2cm}
		
		\consola{========================================================}\par
		
		\vspace{0.2cm}
		
		\consola{Por favor, seleccione una opción (1-5): \textgreater\ }\entrada{1}\par
		
		\vspace{2.8cm}
		
		\lineaPSeInt{línea 41 instrucción 1}
		
	\end{PSeIntWindow}
	%======================TRANSFORMACIO UNIDADES===============================
	\vspace{0.2cm}
	\InicioEjecucion
	{TRANSFORMACION DE UNIDADES}
	{ejec:04-transformacion-de-unidades}
	
	\begin{PSeIntWindow}
		
		\BarraPSeInt
		
		\vspace{0.2cm}
		
		\ttfamily
		\small
		
		\consola{==============================================================}\par
		
		\begin{center}
			\textcolor{pseString}{MODULO 1 \hspace{1cm} TRANSFORMACION DE UNIDADES}
		\end{center}
		
		\consola{==============================================================}\par
		
		\vspace{0.6cm}
		
		\consola{Ingrese la magnitud en Metros a convertir: \textgreater}\par
		
		\vspace{1.8cm}
		
		\lineaPSeInt{línea 58 instrucción 1}
		
	\end{PSeIntWindow}
	
	%=====================EJECUCION 1 DE TRANSFORMACION DE UNIDADES===
	
	\InicioEjecucion
	{EJECUCION 1 TRANSFORMACION DE UNIDADES}
	{ejec:05-ejecucion-1-transformacion-de-unidades}
	
	\begin{PSeIntWindow}
		
		\BarraPSeInt
		
		\vspace{0.2cm}
		
		\ttfamily
		\small
		
		\consola{========================================================}\par
		
		\begin{center}
			\textcolor{pseString}{MODULO 1 \hspace{1cm} TRANSFORMACION DE UNIDADES}
		\end{center}
		
		\consola{========================================================}\par
		
		\vspace{0.4cm}
		
		\consola{Ingrese la magnitud en Metros a convertir: \textgreater\ }\entrada{10}\par
		
		\vspace{0.4cm}
		
		\consola{--------------------------------------------------------}\par
		
		\vspace{0.2cm}
		
		\consola{RESULTADOS DE LA CONVERSIÓN :}\par
		
		\vspace{0.2cm}
		
		\consola{\hspace*{0.5cm}-> Metros origen: }\entrada{10}\consola{ METROS}\par
		
		\consola{\hspace*{0.5cm}-> Pies : }\entrada{32.8084}\consola{ PIES}\par
		
		\consola{\hspace*{0.5cm}-> Pulgadas : }\entrada{393.701}\consola{ PULGADAS}\par
		
		\consola{\hspace*{0.5cm}-> Yardas : }\entrada{10.9361}\consola{ YARDAS}\par
		
		\consola{\hspace*{0.5cm}-> Millas : }\entrada{0.00621371}\consola{ MILLAS}\par
		
		\vspace{0.2cm}
		
		\consola{========================================================}\par
		
		\vspace{0.4cm}
		
		\consola{Presione una tecla para retornar al Menú Principal...}\par
		
		\vspace{1.5cm}
		
		\lineaPSeInt{línea 78 instrucción 1}
		
	\end{PSeIntWindow}
	%===================FIN OPCION 1 DEL MENU =========================
	
	%=========OPCION 1========================
	\vspace{0.2cm}
	\InicioEjecucion
	{EJECUCION 2 MENÚ DE OPCIONES OPCION 1}
	{ejec:06-ejecucion-2-menu-de-opciones-opcion-1}
	
	\begin{PSeIntWindow}
		
		\BarraPSeInt
		
		\vspace{0.2cm}
		
		\ttfamily
		\small
		
		\consola{========================================================}\par
		
		\begin{center}
			\textcolor{pseString}{MENÚ PRINCIPAL}
		\end{center}
		
		\consola{========================================================}\par
		
		\vspace{0.2cm}
		
		\consola{1. Conversión de Longitudes}\par
		
		\consola{2. Calcular el sueldo neto semanal de un trabajador}\par
		
		\consola{3. Suma y resta de dos arreglos unidimensionales}\par
		
		\consola{4. Inversa de un arreglo de N elementos}\par
		
		\consola{5. Salir}\par
		
		\vspace{0.2cm}
		
		\consola{========================================================}\par
		
		\vspace{0.2cm}
		
		\consola{Por favor, seleccione una opción (1-5): \textgreater\ }\entrada{1}\par
		
		\vspace{2.8cm}
		
		\lineaPSeInt{línea 41 instrucción 1}
		
	\end{PSeIntWindow}
	
	%=============================EJECUCION 2 TRANSFORMACION DE UNIDADES
	\vspace{0.2cm}
	\InicioEjecucion
	{EJECUCION 2 TRANSFORMACION DE UNIDADES}
	{ejec:07-ejecucion-2-transformacion-de-unidades}
	
	\begin{PSeIntWindow}
		
		\BarraPSeInt
		
		\vspace{0.2cm}
		
		\ttfamily
		\small
		
		\consola{==============================================================}\par
		
		\begin{center}
			\textcolor{pseString}{MODULO 1 \hspace{1cm} TRANSFORMACION DE UNIDADES}
		\end{center}
		
		\consola{==============================================================}\par
		
		\vspace{0.6cm}
		
		\consola{Ingrese la magnitud en Metros a convertir: \textgreater}\par
		
		\vspace{1.8cm}
		
		\lineaPSeInt{línea 58 instrucción 1}
		
	\end{PSeIntWindow}
	
	%=====================EJECUCION 1 DE tRANSFORMACION DE UNIDADES===
	\vspace{0.2cm}
	\InicioEjecucion
	{EJECUCION 2 TRANSFORMACION DE UNIDADES}
	{ejec:08-ejecucion-2-transformacion-de-unidades}
	
	\begin{PSeIntWindow}
		
		\BarraPSeInt
		
		\vspace{0.2cm}
		
		\ttfamily
		\small
		
		\consola{========================================================}\par
		
		\begin{center}
			\textcolor{pseString}{MODULO 1 \hspace{1cm} TRANSFORMACION DE UNIDADES}
		\end{center}
		
		\consola{========================================================}\par
		
		\vspace{0.4cm}
		
		\consola{Ingrese la magnitud en Metros a convertir: \textgreater\ }\entrada{20}\par
		
		\vspace{0.4cm}
		
		\consola{--------------------------------------------------------}\par
		
		\vspace{0.2cm}
		
		\consola{RESULTADOS DE LA CONVERSIÓN :}\par
		
		\vspace{0.2cm}
		
		\consola{\hspace*{0.5cm}-> Metros origen: }\entrada{20}\consola{ METROS}\par
		
		\consola{\hspace*{0.5cm}-> Pies : }\entrada{65.616}\consola{ PIES}\par
		
		\consola{\hspace*{0.5cm}-> Pulgadas : }\entrada{787.40}\consola{ PULGADAS}\par
		
		\consola{\hspace*{0.5cm}-> Yardas : }\entrada{21.872}\consola{ YARDAS}\par
		
		\consola{\hspace*{0.5cm}-> Millas : }\entrada{0.0124274}\consola{ MILLAS}\par
		
		\vspace{0.2cm}
		
		\consola{========================================================}\par
		
		\vspace{0.4cm}
		
		\consola{Presione una tecla para retornar al Menú Principal...}\par
		
		\vspace{1.5cm}
		
		\lineaPSeInt{línea 78 instrucción 1}
		
	\end{PSeIntWindow}	
	
	%===================EJECUCION MENU PRONCIPAL===================================
	\InicioEjecucion
	{EJECUCION MENU PRINCIPAL}
	{ejec:09-ejecucion-menu-principal}
	
	\begin{PSeIntWindow}
		
		\BarraPSeInt
		
		\vspace{0.2cm}
		
		\ttfamily
		\small
		
		\consola{========================================================}\par
		
		\begin{center}
			\textcolor{pseString}{MENÚ PRINCIPAL}
		\end{center}
		
		\consola{========================================================}\par
		
		\vspace{0.2cm}
		
		\consola{1. Conversión de Longitudes}\par
		
		\consola{2. Calcular el sueldo neto semanal de un trabajador}\par
		
		\consola{3. Suma y resta de dos arreglos unidimensionales}\par
		
		\consola{4. Inversa de un arreglo de N elementos}\par
		
		\consola{5. Salir}\par
		
		\vspace{0.2cm}
		
		\consola{========================================================}\par
		
		\vspace{0.2cm}
		
		\consola{Por favor, seleccione una opción (1-5): \textgreater\ }\entrada{2}\par
		
		\vspace{2.8cm}
		
		\lineaPSeInt{línea 41 instrucción 1}
		
	\end{PSeIntWindow}
	
	%========EJECUCION EJECUCION IGRESO DE DATOS ROL=========
	\vspace{0.2cm}
	\InicioEjecucion
	{EJECUCION IGRESO DE DATOS ROL}
	{ejec:10-ejecucion-igreso-de-datos-rol}
	
	\begin{PSeIntWindow}
		
		\BarraPSeInt
		
		\vspace{0.2cm}
		
		\ttfamily
		\small
		
		\consola{========================================================}\par
		
		\begin{center}
			\textcolor{pseString}{MODULO 2: CALCULO NÓMINA SEMANAL}
		\end{center}
		
		\consola{========================================================}\par
		
		\vspace{0.4cm}
		
		\consola{Ingrese el número total de horas trabajadas en la semana: \textgreater}\par
		
		\vspace{2.8cm}
		
		\lineaPSeInt{línea 102 instrucción 1}
		
	\end{PSeIntWindow}	
	%==========CALCULO DEL ROL======================
	
	\InicioEjecucion
	{EJECUCION DE CALCULO ROL}
	{ejec:11-ejecucion-de-claculo-rol}
	
	\begin{PSeIntWindow}
		
		\BarraPSeInt
		
		\vspace{0.2cm}
		
		\ttfamily
		\small
		
		\consola{========================================================}\par
		
		\begin{center}
			\textcolor{pseString}{MODULO 2: CALCULO NÓMINA SEMANAL}
		\end{center}
		
		\consola{========================================================}\par
		
		\vspace{0.3cm}
		
		\consola{Ingrese el número total de horas trabajadas en la semana: \textgreater\ }\entrada{50}\par
		
		\consola{Ingrese el costo de la hora laboral estándar (Dolares): \textgreater\ }\entrada{4}\par
		
		\vspace{0.3cm}
		
		\consola{--------------------------------------------------------}\par
		
		\vspace{0.2cm}
		
		\consola{ROL DE PAGO:}\par
		
		\vspace{0.2cm}
		
		\consola{\hspace*{0.5cm}[+] Horas Ordinarias Liquidadas: }\entrada{40}\consola{ hrs}\par
		
		\consola{\hspace*{0.5cm}[+] Horas Extraordinarias Liquidadas: }\entrada{10}\consola{ hrs}\par
		
		\consola{\hspace*{0.5cm}[=] Sueldo Bruto Consolidado: \$}\entrada{220}\par
		
		\consola{\hspace*{0.5cm}[-] Descuento Obligatorio IESS (9.2\%): \$}\entrada{20.24}\par
		
		\consola{\hspace*{0.5cm}[TOTAL A PAGAR] Sueldo Neto Semanal: \$}\entrada{199.76}\par
		
		\vspace{0.2cm}
		
		\consola{========================================================}\par
		
		\vspace{0.4cm}
		
		\consola{Presione una tecla para retornar al Menú Principal...}\par
		
		\vspace{1.8cm}
		
		\lineaPSeInt{línea 125 instrucción 1}
		
	\end{PSeIntWindow}	
	%======================SEGUNDA EJECUCION=================
	\vspace{0.2cm}
	\InicioEjecucion
	{EJECUCION 2 CALCULO DE ROL}
	{ejec:12-ejecucion-2-claculo-de-rol}
	
	\begin{PSeIntWindow}
		
		\BarraPSeInt
		
		\vspace{0.2cm}
		
		\ttfamily
		\small
		
		\consola{========================================================}\par
		
		\begin{center}
			\textcolor{pseString}{MODULO 2: CALCULO NÓMINA SEMANAL}
		\end{center}
		
		\consola{========================================================}\par
		
		\vspace{0.3cm}
		
		\consola{Ingrese el número total de horas trabajadas en la semana: \textgreater\ }\entrada{40}\par
		
		\consola{Ingrese el costo de la hora laboral estándar (Dolares): \textgreater\ }\entrada{4}\par
		
		\vspace{0.3cm}
		
		\consola{--------------------------------------------------------}\par
		
		\vspace{0.2cm}
		
		\consola{ROL DE PAGO:}\par
		
		\vspace{0.2cm}
		
		\consola{\hspace*{0.5cm}[+] Horas Ordinarias Liquidadas: }\entrada{40}\consola{ hrs}\par
		
		\consola{\hspace*{0.5cm}[+] Horas Extraordinarias Liquidadas: }\entrada{0}\consola{ hrs}\par
		
		\consola{\hspace*{0.5cm}[=] Sueldo Bruto Consolidado: \$}\entrada{160}\par
		
		\consola{\hspace*{0.5cm}[-] Descuento Obligatorio IESS (9.2\%): \$}\entrada{14.72}\par
		
		\consola{\hspace*{0.5cm}[TOTAL A PAGAR] Sueldo Neto Semanal: \$}\entrada{145.28}\par
		
		\vspace{0.2cm}
		
		\consola{========================================================}\par
		
		\vspace{0.4cm}
		
		\consola{Presione una tecla para retornar al Menú Principal...}\par
		
		\vspace{1.8cm}
		
		\lineaPSeInt{línea 122 instrucción 1}
		
	\end{PSeIntWindow}	
	%======================OPCION 3 DEL MENU ===================
	\vspace{0.2cm}
	\InicioEjecucion
	{EJECUCION MENU PRINCIPAL}
	{ejec:13-ejecucion-menu-principal}
	
	\begin{PSeIntWindow}
		
		\BarraPSeInt
		
		\vspace{0.2cm}
		
		\ttfamily
		\small
		
		\consola{========================================================}\par
		
		\begin{center}
			\textcolor{pseString}{MENÚ PRINCIPAL}
		\end{center}
		
		\consola{========================================================}\par
		
		\vspace{0.2cm}
		
		\consola{1. Conversión de Longitudes}\par
		
		\consola{2. Calcular el sueldo neto semanal de un trabajador}\par
		
		\consola{3. Suma y resta de dos arreglos unidimensionales}\par
		
		\consola{4. Inversa de un arreglo de N elementos}\par
		
		\consola{5. Salir}\par
		
		\vspace{0.2cm}
		
		\consola{========================================================}\par
		
		\vspace{0.2cm}
		
		\consola{Por favor, seleccione una opción (1-5): \textgreater\ }\entrada{3}\par
		
		\vspace{2.8cm}
		
		\lineaPSeInt{línea 41 instrucción 1}
		
	\end{PSeIntWindow}	
	%==========================modulo 3=============
	\vspace{0.2cm}
	\InicioEjecucion
	{MÓDULO 3 - INGRESO DE DATOS DE VECTORES}
	{ejec:14-modulo-3-ingreso-de-datos-de-vectores}
	
	\begin{PSeIntWindow}
		
		\BarraPSeInt
		
		\vspace{0.2cm}
		
		\ttfamily
		\small
		
		\consola{========================================================}\par
		
		\begin{center}
			\textcolor{pseString}{MODULO 3: OPERACIONES CON VECTORES (SUMA-RESTA)}
		\end{center}
		
		\consola{========================================================}\par
		
		\vspace{0.2cm}
		
		\consola{\textgreater\textgreater\ ESCRIBIR 10 ELEMENTOS PARA EL VECTOR A:}\par
		
		\vspace{0.15cm}
		
		\consola{Ingrese numero VectorA[1]: \textgreater\ }\entrada{4}\par
		\consola{Ingrese numero VectorA[2]: \textgreater\ }\entrada{5}\par
		\consola{Ingrese numero VectorA[3]: \textgreater\ }\entrada{3}\par
		\consola{Ingrese numero VectorA[4]: \textgreater\ }\entrada{2}\par
		\consola{Ingrese numero VectorA[5]: \textgreater\ }\entrada{6}\par
		\consola{Ingrese numero VectorA[6]: \textgreater\ }\entrada{7}\par
		\consola{Ingrese numero VectorA[7]: \textgreater\ }\entrada{8}\par
		\consola{Ingrese numero VectorA[8]: \textgreater\ }\entrada{9}\par
		\consola{Ingrese numero VectorA[9]: \textgreater\ }\entrada{3}\par
		\consola{Ingrese numero VectorA[10]: \textgreater\ }\entrada{2}\par
		
		\vspace{0.2cm}
		
		\consola{--------------------------------------------------------}\par
		
		\vspace{0.2cm}
		
		\consola{\textgreater\textgreater\ ESCRIBIR 10 ELEMENTOS PARA EL VECTOR B:}\par
		
		\vspace{0.15cm}
		
		\consola{Ingrese numero VectorB[1]: \textgreater\ }\entrada{1}\par
		\consola{Ingrese numero VectorB[2]: \textgreater\ }\entrada{2}\par
		\consola{Ingrese numero VectorB[3]: \textgreater\ }\entrada{3}\par
		\consola{Ingrese numero VectorB[4]: \textgreater\ }\entrada{4}\par
		\consola{Ingrese numero VectorB[5]: \textgreater\ }\entrada{5}\par
		\consola{Ingrese numero VectorB[6]: \textgreater\ }\entrada{6}\par
		\consola{Ingrese numero VectorB[7]: \textgreater\ }\entrada{7}\par
		\consola{Ingrese numero VectorB[8]: \textgreater\ }\entrada{8}\par
		\consola{Ingrese numero VectorB[9]: \textgreater\ }\entrada{9}\par
		\consola{Ingrese numero VectorB[10]: \textgreater\ }\entrada{10}\par
		
		\vspace{0.5cm}
		
		\lineaPSeInt{línea 226 instrucción 1}
		
	\end{PSeIntWindow}
	
	%===============================resultado de suma y rsta de vectores 1ejecucion
	\vspace{0.2cm}
	\InicioEjecucion
	{RESULTADOS DE SUMA Y RESTA DE VECTORES}
	{ejec:15-resultados-de-suma-y-resta-de-vectores}
	
	\begin{PSeIntWindow}
		
		\BarraPSeInt
		
		\vspace{0.2cm}
		
		\ttfamily\small
		
		\consola{==============================================================}\par
		
		\begin{center}
			\textcolor{pseString}{IMPRESION DE VALORES ORDENADOS}
		\end{center}
		
		\consola{==============================================================}\par
		
		\vspace{0.25cm}
		
		\consola{Índice \hspace{0.2cm}| \hspace{0.2cm} Vector A \hspace{0.2cm}| \hspace{0.2cm} Vector B \hspace{0.2cm}| \hspace{0.2cm} Suma (A + B) \hspace{0.2cm}| \hspace{0.2cm} Resta (A - B)}\par
		\consola{-------------------------------------------------------------------------------------------------------------------------------}\par
		
		\consola{[1]  \hspace{0.7cm}| \hspace{0.9cm}}\entrada{4}\consola{\hspace{1.0cm}| \hspace{0.9cm}}\entrada{1}\consola{\hspace{1.0cm}| \hspace{1.0cm}}\entrada{5}\consola{\hspace{1.5cm}| \hspace{1.3cm}}\entrada{3}\par
		\consola{[2]  \hspace{0.7cm}| \hspace{0.9cm}}\entrada{5}\consola{\hspace{1.0cm}| \hspace{0.9cm}}\entrada{2}\consola{\hspace{1.0cm}| \hspace{1.0cm}}\entrada{7}\consola{\hspace{1.5cm}| \hspace{1.3cm}}\entrada{3}\par
		\consola{[3]  \hspace{0.7cm}| \hspace{0.9cm}}\entrada{3}\consola{\hspace{1.0cm}| \hspace{0.9cm}}\entrada{3}\consola{\hspace{1.0cm}| \hspace{1.0cm}}\entrada{6}\consola{\hspace{1.5cm}| \hspace{1.3cm}}\entrada{0}\par
		\consola{[4]  \hspace{0.7cm}| \hspace{0.9cm}}\entrada{2}\consola{\hspace{1.0cm}| \hspace{0.9cm}}\entrada{4}\consola{\hspace{1.0cm}| \hspace{1.0cm}}\entrada{6}\consola{\hspace{1.5cm}| \hspace{1.1cm}}\entrada{-2}\par
		\consola{[5]  \hspace{0.7cm}| \hspace{0.9cm}}\entrada{6}\consola{\hspace{1.0cm}| \hspace{0.9cm}}\entrada{5}\consola{\hspace{1.0cm}| \hspace{0.8cm}}\entrada{11}\consola{\hspace{1.4cm}| \hspace{1.3cm}}\entrada{1}\par
		\consola{[6]  \hspace{0.7cm}| \hspace{0.9cm}}\entrada{7}\consola{\hspace{1.0cm}| \hspace{0.9cm}}\entrada{6}\consola{\hspace{1.0cm}| \hspace{0.8cm}}\entrada{13}\consola{\hspace{1.4cm}| \hspace{1.3cm}}\entrada{1}\par
		\consola{[7]  \hspace{0.7cm}| \hspace{0.9cm}}\entrada{8}\consola{\hspace{1.0cm}| \hspace{0.9cm}}\entrada{7}\consola{\hspace{1.0cm}| \hspace{0.8cm}}\entrada{15}\consola{\hspace{1.4cm}| \hspace{1.3cm}}\entrada{1}\par
		\consola{[8]  \hspace{0.7cm}| \hspace{0.9cm}}\entrada{9}\consola{\hspace{1.0cm}| \hspace{0.9cm}}\entrada{8}\consola{\hspace{1.0cm}| \hspace{0.8cm}}\entrada{17}\consola{\hspace{1.4cm}| \hspace{1.3cm}}\entrada{1}\par
		\consola{[9]  \hspace{0.7cm}| \hspace{0.9cm}}\entrada{3}\consola{\hspace{1.0cm}| \hspace{0.9cm}}\entrada{9}\consola{\hspace{1.0cm}| \hspace{0.8cm}}\entrada{12}\consola{\hspace{1.4cm}| \hspace{1.1cm}}\entrada{-6}\par
		\consola{[10] \hspace{0.5cm}| \hspace{0.9cm}}\entrada{2}\consola{\hspace{1.0cm}| \hspace{0.7cm}}\entrada{10}\consola{\hspace{0.8cm}| \hspace{0.8cm}}\entrada{12}\consola{\hspace{1.4cm}| \hspace{1.1cm}}\entrada{-8}\par
		
		\vspace{0.25cm}
		
		\consola{==============================================================}\par
		
		\vspace{0.25cm}
		
		\consola{Presione una tecla para retornar al Menú Principal...}\par
		
		\vspace{0.4cm}
		
		\lineaPSeInt{línea 245 instrucción 1}
		
	\end{PSeIntWindow}
	
	%==========================ejecucion 2==============================
	\vspace{0.2cm}
	\InicioEjecucion
	{EJECUCION 2 OPERACIONES CON VECTORES (SUMA-RESTA)}
	{ejec:16-ejecucion-2-operaciones-con-vectores-suma-resta}
	
	\begin{PSeIntWindow}
		
		\BarraPSeInt
		
		\vspace{0.2cm}
		
		\ttfamily
		\small
		
		\consola{========================================================}\par
		
		\begin{center}
			\textcolor{pseString}{MODULO 3: OPERACIONES CON VECTORES (SUMA-RESTA)}
		\end{center}
		
		\consola{========================================================}\par
		
		\vspace{0.2cm}
		
		\consola{\textgreater\textgreater\ ESCRIBIR 10 ELEMENTOS PARA EL VECTOR A:}\par
		
		\vspace{0.15cm}
		
		\consola{Ingrese numero VectorA[1]: \textgreater\ }\entrada{6}\par
		\consola{Ingrese numero VectorA[2]: \textgreater\ }\entrada{7}\par
		\consola{Ingrese numero VectorA[3]: \textgreater\ }\entrada{8}\par
		\consola{Ingrese numero VectorA[4]: \textgreater\ }\entrada{6}\par
		\consola{Ingrese numero VectorA[5]: \textgreater\ }\entrada{56}\par
		\consola{Ingrese numero VectorA[6]: \textgreater\ }\entrada{6}\par
		\consola{Ingrese numero VectorA[7]: \textgreater\ }\entrada{7}\par
		\consola{Ingrese numero VectorA[8]: \textgreater\ }\entrada{8}\par
		\consola{Ingrese numero VectorA[9]: \textgreater\ }\entrada{5}\par
		\consola{Ingrese numero VectorA[10]: \textgreater\ }\entrada{4}\par
		
		\vspace{0.2cm}
		
		\consola{--------------------------------------------------------}\par
		
		\vspace{0.2cm}
		
		\consola{\textgreater\textgreater\ ESCRIBIR 10 ELEMENTOS PARA EL VECTOR B:}\par
		
		\vspace{0.15cm}
		
		\consola{Ingrese numero VectorB[1]: \textgreater\ }\entrada{9}\par
		\consola{Ingrese numero VectorB[2]: \textgreater\ }\entrada{7}\par
		\consola{Ingrese numero VectorB[3]: \textgreater\ }\entrada{6}\par
		\consola{Ingrese numero VectorB[4]: \textgreater\ }\entrada{5}\par
		\consola{Ingrese numero VectorB[5]: \textgreater\ }\entrada{4}\par
		\consola{Ingrese numero VectorB[6]: \textgreater\ }\entrada{43}\par
		\consola{Ingrese numero VectorB[7]: \textgreater\ }\entrada{3}\par
		\consola{Ingrese numero VectorB[8]: \textgreater\ }\entrada{3}\par
		\consola{Ingrese numero VectorB[9]: \textgreater\ }\entrada{3}\par
		\consola{Ingrese numero VectorB[10]: \textgreater\ }\entrada{6}\par
		
		\vspace{0.5cm}
		
		\lineaPSeInt{línea 226 instrucción 1}
		
	\end{PSeIntWindow}
	
	%==================================ingreso vector=================
	
	\InicioEjecucion
	{EJECUCION 2 RESULTADOS DE SUMA Y RESTA DE VECTORES}
	{ejec:17-ejecucion-2-resultados-de-suma-y-resta-de-vectores}
	
	\begin{PSeIntWindow}
		
		\BarraPSeInt
		
		\vspace{0.2cm}
		
		\ttfamily\small
		
		\consola{========================================================}\par
		
		\begin{center}
			\textcolor{pseString}{IMPRESION DE VALORES ORDENADOS}
		\end{center}
		
		\consola{========================================================}\par
		
		\vspace{0.25cm}
		
		\consola{Índice \hspace{0.4cm}| \hspace{0.4cm} Vector A \hspace{0.4cm}| \hspace{0.4cm} Vector B \hspace{0.4cm}| \hspace{0.4cm} Suma (A + B) \hspace{0.4cm}| \hspace{0.4cm} Resta (A - B)}\par
		\consola{----------------------------------------------------------------}\par
		
		\consola{[1]  \hspace{0.7cm}| \hspace{0.9cm}}\entrada{6}\consola{\hspace{1.0cm}| \hspace{0.9cm}}\entrada{9}\consola{\hspace{1.0cm}| \hspace{0.8cm}}\entrada{15}\consola{\hspace{1.4cm}| \hspace{1.1cm}}\entrada{-3}\par
		\consola{[2]  \hspace{0.7cm}| \hspace{0.9cm}}\entrada{7}\consola{\hspace{1.0cm}| \hspace{0.9cm}}\entrada{7}\consola{\hspace{1.0cm}| \hspace{0.8cm}}\entrada{14}\consola{\hspace{1.4cm}| \hspace{1.3cm}}\entrada{0}\par
		\consola{[3]  \hspace{0.7cm}| \hspace{0.9cm}}\entrada{8}\consola{\hspace{1.0cm}| \hspace{0.9cm}}\entrada{6}\consola{\hspace{1.0cm}| \hspace{0.8cm}}\entrada{14}\consola{\hspace{1.4cm}| \hspace{1.3cm}}\entrada{2}\par
		\consola{[4]  \hspace{0.7cm}| \hspace{0.9cm}}\entrada{6}\consola{\hspace{1.0cm}| \hspace{0.9cm}}\entrada{5}\consola{\hspace{1.0cm}| \hspace{0.8cm}}\entrada{11}\consola{\hspace{1.4cm}| \hspace{1.3cm}}\entrada{1}\par
		\consola{[5]  \hspace{0.7cm}| \hspace{0.6cm}}\entrada{56}\consola{\hspace{0.8cm}| \hspace{0.9cm}}\entrada{4}\consola{\hspace{1.0cm}| \hspace{0.8cm}}\entrada{60}\consola{\hspace{1.4cm}| \hspace{1.1cm}}\entrada{52}\par
		\consola{[6]  \hspace{0.7cm}| \hspace{0.9cm}}\entrada{6}\consola{\hspace{1.0cm}| \hspace{0.6cm}}\entrada{43}\consola{\hspace{0.8cm}| \hspace{0.8cm}}\entrada{49}\consola{\hspace{1.4cm}| \hspace{1.0cm}}\entrada{-37}\par
		\consola{[7]  \hspace{0.7cm}| \hspace{0.9cm}}\entrada{7}\consola{\hspace{1.0cm}| \hspace{0.9cm}}\entrada{3}\consola{\hspace{1.0cm}| \hspace{0.8cm}}\entrada{10}\consola{\hspace{1.4cm}| \hspace{1.3cm}}\entrada{4}\par
		\consola{[8]  \hspace{0.7cm}| \hspace{0.9cm}}\entrada{8}\consola{\hspace{1.0cm}| \hspace{0.9cm}}\entrada{3}\consola{\hspace{1.0cm}| \hspace{0.8cm}}\entrada{11}\consola{\hspace{1.4cm}| \hspace{1.3cm}}\entrada{5}\par
		\consola{[9]  \hspace{0.7cm}| \hspace{0.9cm}}\entrada{5}\consola{\hspace{1.0cm}| \hspace{0.9cm}}\entrada{3}\consola{\hspace{1.0cm}| \hspace{1.0cm}}\entrada{8}\consola{\hspace{1.5cm}| \hspace{1.3cm}}\entrada{2}\par
		\consola{[10] \hspace{0.5cm}| \hspace{0.9cm}}\entrada{4}\consola{\hspace{1.0cm}| \hspace{0.9cm}}\entrada{6}\consola{\hspace{1.0cm}| \hspace{0.8cm}}\entrada{10}\consola{\hspace{1.4cm}| \hspace{1.1cm}}\entrada{-2}\par
		
		\vspace{0.25cm}
		
		\consola{========================================================}\par
		
		\vspace{0.25cm}
		
		\consola{Presione una tecla para retornar al Menú Principal...}\par
		
		\vspace{0.4cm}
		
		\lineaPSeInt{línea 245 instrucción 1}
		
	\end{PSeIntWindow}
	%	=======================opcion 4================================
	\vspace{0.2cm}
	\InicioEjecucion
	{INGRESO DEL TAMAÑO DEL VECTOR}
	{ejec:18-ingreso-del-tamano-del-vector}
	
	\begin{PSeIntWindow}
		
		\BarraPSeInt
		
		\vspace{0.2cm}
		
		\ttfamily\small
		
		\consola{========================================================}\par
		
		\begin{center}
			\textcolor{pseString}{MODULO 4: INVERSO DEL VECTOR INGRESADO}
		\end{center}
		
		\consola{========================================================}\par
		
		\vspace{0.4cm}
		
		\consola{Defina el tamaño del vector debe ser positivo: \textgreater\ }\par
		
		\vspace{2.5cm}
		
		\lineaPSeInt{línea 273 instrucción 1}
		
	\end{PSeIntWindow}
	
	%=================resultado moduo 4====
	
	\InicioEjecucion
	{INGRESO DE DATOS DEL VECTOR}
	{ejec:19-ingreso-de-datos-del-vector}
	
	\begin{PSeIntWindow}
		
		\BarraPSeInt
		
		\vspace{0.2cm}
		
		\ttfamily\small
		
		\consola{========================================================}\par
		
		\begin{center}
			\textcolor{pseString}{MODULO 4: INVERSO DEL VECTOR INGRESADO}
		\end{center}
		
		\consola{========================================================}\par
		
		\vspace{0.25cm}
		
		\consola{Defina el tamaño del vector debe ser positivo: \textgreater\ }\entrada{9}\par
		
		\consola{Ingrese valor para la posicion [1]: \textgreater\ }\entrada{1}\par
		
		\consola{Ingrese valor para la posicion [2]: \textgreater\ }\entrada{2}\par
		
		\consola{Ingrese valor para la posicion [3]: \textgreater\ }\entrada{3}\par
		
		\consola{Ingrese valor para la posicion [4]: \textgreater\ }\entrada{4}\par
		
		\consola{Ingrese valor para la posicion [5]: \textgreater\ }\entrada{5}\par
		
		\consola{Ingrese valor para la posicion [6]: \textgreater\ }\entrada{6}\par
		
		\consola{Ingrese valor para la posicion [7]: \textgreater\ }\entrada{7}\par
		
		\consola{Ingrese valor para la posicion [8]: \textgreater\ }\entrada{8}\par
		
		\consola{Ingrese valor para la posicion [9]: \textgreater\ }\entrada{9}\par
		
		\vspace{1.5cm}
		
		\lineaPSeInt{línea 287 instrucción 1}
		
	\end{PSeIntWindow}
	%==========================resiltado modulo 4 ==============
	\vspace{0.2cm}
	\InicioEjecucion
	{VECTOR ORDENADO E INVERSO}
	{ejec:20-vector-ordenado-e-inverso}
	
	\begin{PSeIntWindow}
		
		\BarraPSeInt
		
		\vspace{0.2cm}
		
		\ttfamily\small
		
		\consola{========================================================}\par
		
		\vspace{0.2cm}
		
		\consola{VECTOR ORDENADO}\par
		
		\vspace{0.2cm}
		
		\consola{========================================================}\par
		
		\vspace{0.3cm}
		
		\consola{[*] Elemento en Celda [1] = }\entrada{1}\par
		
		\consola{[*] Elemento en Celda [2] = }\entrada{2}\par
		
		\consola{[*] Elemento en Celda [3] = }\entrada{3}\par
		
		\consola{[*] Elemento en Celda [4] = }\entrada{4}\par
		
		\consola{[*] Elemento en Celda [5] = }\entrada{5}\par
		
		\consola{[*] Elemento en Celda [6] = }\entrada{6}\par
		
		\consola{[*] Elemento en Celda [7] = }\entrada{7}\par
		
		\consola{[*] Elemento en Celda [8] = }\entrada{8}\par
		
		\consola{[*] Elemento en Celda [9] = }\entrada{9}\par
		
		\vspace{0.3cm}
		
		\consola{========================================================}\par
		
		\vspace{0.3cm}
		
		\consola{ORDEN INVERSO DEL VECTOR}\par
		
		\vspace{0.2cm}
		
		\consola{========================================================}\par
		
		\vspace{0.3cm}
		
		\consola{[\#] Elemento en indice Invertido [9] = }\entrada{9}\par
		
		\consola{[\#] Elemento en indice Invertido [8] = }\entrada{8}\par
		
		\consola{[\#] Elemento en indice Invertido [7] = }\entrada{7}\par
		
		\consola{[\#] Elemento en indice Invertido [6] = }\entrada{6}\par
		
		\consola{[\#] Elemento en indice Invertido [5] = }\entrada{5}\par
		
		\consola{[\#] Elemento en indice Invertido [4] = }\entrada{4}\par
		
		\consola{[\#] Elemento en indice Invertido [3] = }\entrada{3}\par
		
		\consola{[\#] Elemento en indice Invertido [2] = }\entrada{2}\par
		
		\consola{[\#] Elemento en indice Invertido [1] = }\entrada{1}\par
		
		\vspace{0.3cm}
		
		\consola{========================================================}\par
		
		\vspace{0.3cm}
		
		\consola{Presione una tecla para retornar al Menú Principal...}\par
		
		\vspace{0.4cm}
		
		\lineaPSeInt{línea 314 instrucción 1}
		
	\end{PSeIntWindow}
	%==============ingreso de datos ejecucion 2=======
	\vspace{0.2cm}
	\InicioEjecucion
	{EJEUCION 2 INGRESO DE DATOS DEL VECTOR}
	{ejec:21-ejeucion-2-ingreso-de-datos-del-vector}
	
	\begin{PSeIntWindow}
		
		\BarraPSeInt
		
		\vspace{0.2cm}
		
		\ttfamily\small
		
		\consola{========================================================}\par
		
		\begin{center}
			\textcolor{pseString}{MODULO 4: INVERSO DEL VECTOR INGRESADO}
		\end{center}
		
		\consola{========================================================}\par
		
		\vspace{0.25cm}
		
		\consola{Defina el tamaño del vector debe ser positivo: \textgreater\ }\entrada{7}\par
		
		\consola{Ingrese valor para la posicion [1]: \textgreater\ }\entrada{10}\par
		
		\consola{Ingrese valor para la posicion [2]: \textgreater\ }\entrada{20}\par
		
		\consola{Ingrese valor para la posicion [3]: \textgreater\ }\entrada{30}\par
		
		\consola{Ingrese valor para la posicion [4]: \textgreater\ }\entrada{40}\par
		
		\consola{Ingrese valor para la posicion [5]: \textgreater\ }\entrada{50}\par
		
		\consola{Ingrese valor para la posicion [6]: \textgreater\ }\entrada{60}\par
		
		\consola{Ingrese valor para la posicion [7]: \textgreater\ }\entrada{70}\par
		
		\vspace{1.8cm}
		
		\lineaPSeInt{línea 287 instrucción 1}
		
	\end{PSeIntWindow}
	
	%==================resultado de la ejecucion 2================
	\vspace{0.2cm}
	\InicioEjecucion
	{EJECUCION 2 VECTOR ORDENADO E INVERSO}
	{ejec:22-ejecucion-2-vector-ordenado-e-inverso}
	
	\begin{PSeIntWindow}
		
		\BarraPSeInt
		
		\vspace{0.2cm}
		
		\ttfamily\small
		
		\consola{========================================================}\par
		
		\vspace{0.2cm}
		
		\consola{VECTOR ORDENADO}\par
		
		\vspace{0.2cm}
		
		\consola{========================================================}\par
		
		\vspace{0.3cm}
		
		\consola{[*] Elemento en Celda [1] = }\entrada{10}\par
		
		\consola{[*] Elemento en Celda [2] = }\entrada{20}\par
		
		\consola{[*] Elemento en Celda [3] = }\entrada{30}\par
		
		\consola{[*] Elemento en Celda [4] = }\entrada{40}\par
		
		\consola{[*] Elemento en Celda [5] = }\entrada{50}\par
		
		\consola{[*] Elemento en Celda [6] = }\entrada{60}\par
		
		\consola{[*] Elemento en Celda [7] = }\entrada{70}\par
		
		\vspace{0.3cm}
		
		\consola{========================================================}\par
		
		\vspace{0.3cm}
		
		\consola{ORDEN INVERSO DEL VECTOR}\par
		
		\vspace{0.2cm}
		
		\consola{========================================================}\par
		
		\vspace{0.3cm}
		
		\consola{[\#] Elemento en indice Invertido [7] = }\entrada{70}\par
		
		\consola{[\#] Elemento en indice Invertido [6] = }\entrada{60}\par
		
		\consola{[\#] Elemento en indice Invertido [5] = }\entrada{50}\par
		
		\consola{[\#] Elemento en indice Invertido [4] = }\entrada{40}\par
		
		\consola{[\#] Elemento en indice Invertido [3] = }\entrada{30}\par
		
		\consola{[\#] Elemento en indice Invertido [2] = }\entrada{20}\par
		
		\consola{[\#] Elemento en indice Invertido [1] = }\entrada{10}\par
		
		\vspace{0.3cm}
		
		\consola{========================================================}\par
		
		\vspace{0.3cm}
		
		\consola{Presione una tecla para retornar al Menú Principal...}\par
		
		\vspace{0.4cm}
		
		\lineaPSeInt{línea 314 instrucción 1}
		
	\end{PSeIntWindow}
	
	%=====================================================
	
    \newpage
    
	\listadefiguraspseint
	%***********************************************
	%************************************************
	% =====================================================
	% CONCLUSIONES Y RECOMENDACIONES
	% =====================================================
	
	\newpage
	
	\section{\textcolor{azulUNEMI}{Conclusiones}}
	
	\begin{enumerate}
		\item La modularización mediante subprocesos permitió organizar el programa en módulos independientes, facilitando su comprensión, mantenimiento y reutilización.
		\item El uso de estructuras repetitivas como Para y Repetir permitió automatizar procesos de ingreso, recorrido y presentación de datos, reduciendo la cantidad de instrucciones necesarias.
		\item La implementación de arreglos unidimensionales fortaleció la comprensión sobre el almacenamiento y manipulación de conjuntos de datos dentro de un programa.
		\item La integración de varios ejercicios en un único sistema demostró la importancia de planificar adecuadamente la estructura de un programa para optimizar recursos y mejorar la experiencia del usuario.
		\item El desarrollo de este proyecto contribuyó al fortalecimiento de habilidades fundamentales en lógica de programación, análisis de problemas y diseño de algoritmos, bases necesarias para asignaturas más avanzadas de la carrera.
	\end{enumerate}
	
	\section{\textcolor{azulUNEMI}{Recomendaciones}}
	
	
	\begin{enumerate}
		\item Mantener una adecuada modularización del código para facilitar futuras modificaciones y ampliaciones del sistema.
		\item Implementar validaciones adicionales que controlen tipos de datos incorrectos y eviten errores durante la ejecución.
		\item Documentar cada subproceso mediante comentarios descriptivos para mejorar la comprensión del código por parte de otros desarrolladores.
		\item Realizar pruebas con diferentes casos de entrada, incluyendo datos válidos e inválidos, para verificar el correcto funcionamiento de cada módulo.
		\item Continuar practicando el uso de arreglos, estructuras repetitivas y funciones, ya que constituyen elementos fundamentales en el desarrollo de software.
	\end{enumerate}
	
	
	% =====================================================
	% BIBLIOGRAFÍA EN ÚLTIMA HOJA
	% =====================================================
	
	\newpage
	
	\section{\textcolor{azulUNEMI}{Bibliografía}}
	
	
	\begin{itemize}
		\item Joyanes Aguilar, L. (2020). \textit{Fundamentos de programación: Algoritmos, estructuras de datos y objetos} (5.ª ed.). McGraw-Hill.
		
		\item Pérez López, C. (2021). \textit{Programación estructurada y algoritmos}. Alfaomega.
		
		\item PSeInt. (2025). \textit{PSeInt: Herramienta para asistir a estudiantes en programación}. https://pseint.sourceforge.net
		
		\item Deitel, P., \& Deitel, H. (2020). \textit{Cómo programar} (10.ª ed.). Pearson Educación.
	\end{itemize}
	%************************************************
	%************************************************
	\newpage
	
	\tableofcontents
	
\end{document}
